\chapter*{Conclusion}
\label{ch:conclusion}
\addcontentsline{toc}{chapter}{\nameref{ch:conclusion}}
\section*{Conclusion}
\chaptermark{Conclusion}

In this thesis we classified the regular Lagrangian torus fibrations over
compact surfaces up to Lagrangian equivalence. In order to do this, we showed
that there are certain invariants that we can associate to such a fibration,
and that these invariants completely determine the nature of the fibration.

The first of these invariants is an integral affine structure on the base space
$B$ of the fibration. This completely determines the local nature of any
Lagrangian fibration through the symplectic reference fibration. In Chapter
\ref{ch:nilpotent-hol} we develop the tools required to classify the integral
affine structures on compact surfaces. This turns out to depend on group
theoretic properties of the fundamental group of the surface. The main result
of this chapter is that completeness of an affine manifold is equivalent to
parallel volume if the fundamental group is nilpotent.

This result reduces the problem of determining the integral affine structures
on a compact surface to a problem in simple group theory and linear algebra. In
Chapter \ref{ch:class-IAS}, we explicitly compute all of the integral affine
structures on the two torus and the Klein bottle. The integral affine
structures on the two torus are classified in Theorem
\ref{thm:classification-torus}. These were found to be the following.

\classificationtorus*

Two integral affine structures from distinct families are never equivalent.
Integral affine structures in the family $\mathcal{A}_{a,b,c,d}$ are sometimes
equivalent as illustrated in Theorem \ref{thm:torus-series-a}.

\torusA*
This completes the full classification of the integral affine structures on the two torus.

We also classify the integral affine structures on the Klein bottle. This is
the result of Theorem \ref{thm:IAS-Klein}. We find that there are four distinct
families and that integral affine structures from distinct families are never
equivalent.

\IASKlein*

The integral affine structure on the base is one invariant of the Lagrangian
fibration. However, it is not the only one. The contents of Chapter
\ref{LagrangianFibrations} and Chapter \ref{ch:classify-fibr} show that there
are two more invariants one can associate to such a fibration. These are the
\emph{Chern class} and the \emph{Lagrangian class}. The Chern class determines
the topological nature of the fibration while the Lagrangian class is a
refinement of this and also controls the symplectic structure. The Lagrangian
class always maps onto the Chern class in the long exact sequence in cohomology
associated to the following short exact sequence in sheaf cohomology;
\begin{align*}
    0 \to \Lambda \to \mathcal{Z}(T^*B) \to \mathcal{Z}(T^*B/\Lambda) \to 0
\end{align*}

where $\Lambda$ is the sheaf of sections of the lattice in $T^*B$ inducing the
integral affine structure and $\mathcal{Z}(\cdot)$ denotes the Lagrangian or
closed sections of the corresponding bundles.

Once we have fixed a Chern class we can ask which Lagrangian classes map onto
it. This corresponds to fixing a topological type for the fibration and asking
which fibrations are topologically equivalent but not necessarily
symplectically. The space of Lagrangian classes which map onto a fix Chern
class turns out to be a quotient of the second de Rham cohomology group. For
the Klein bottle, the second cohomology class vanishes and hence for a fixed
Chern class, all of the Lagrangian fibrations over the Klein bottle which
induce this Chern class are symplectically equivalent. For the torus this is
not the case; there are Lagrangian fibrations which are topologically
equivalent but not symplectically. In both cases we compute the possible Chern
classes in Chapter \ref{ch:classify-fibr}. For the torus this is the content of
Lemma \ref{lemma:chern-class-A} and Lemma \ref{lemma:chern-class-B} which we
summarize below.

\chernTorus*

Lagrangian fibrations over the torus can be topologically equivalent as bundles
but not symplectically equivalent. This is governed by cohomological invariants
of shifts which are the differentials of latticed one-forms. These form a
subgroup of the second De Rham cohomology group and this subgroup parametrizes
Lagrangian classes which map to a fixed Chern class. The calculation of which
classes in $H^2(T^2, \R)$ correspond to cohomological invariants of shifts is
the content of Theorem \ref{thm:invar-calculation}.

For the Klein bottle, every Lagrangian torus fibration which is topologically
equivalent is also symplectically equivalent. The possible Chern classes for
the Klein bottle are found in Theorem \ref{thm:Klein-chern}.

\kleinChern*

In order to explicitly construct a Lagrangian fibration with a fixed Chern
class, we can perform Luttinger surgery as described in Section \ref{sec:Lutt}.
For the Klein bottle, this completes the construction. For the torus we must
add a two form to the base considered up to a cohomological invariant of a
shift. This is described in Section \ref{sec:construct}.

\section*{Whats next?}

There are many directions one can go from here. We will now outline some
possible extensions of the work done in this thesis

\subsection*{Integral affine $3$-manifolds}
In \cite{kozlov2018integralaffine3manifolds}, a classification of complete
integral affine manifolds on compact three-dimensional manifolds up to a
finite-sheeted covering is given. One could attempt to use the methods
developed in this thesis in order to compute the possible Chern and Lagrangian
classes for Lagrangian fibrations over these manifolds. In this paper they find
that there are essentially three different families of such manifold. Up to
finite index, the first families has abelian fundamental group, the second
family has nilpotent fundamental group and the third class has solvable
fundamental group.

For example, the third family which has the most complicated fundamental group,
are mapping tori of the torus by some diffeomrophism $A: T^2 \to T^2$, at least
up to a finite covering. The fundamental group of a mapping torus sits into a
short exact sequence $$ 1 \to \Z^2 \to \Z^2 \rtimes_{A} \Z \to \Z \to 1 $$ and
so in principal, one could try to apply the spectral sequence methods we
developed inSection \ref{sec:spectral-seq} to compute the possible Chern
classes for such manifolds.

\subsection*{Integral-integral affine manifolds}
One could further try to refine the classification of integral affine manifolds
to the classification of integral-integral affine manifolds over compact
surfaces. An integral-integral affine manifold is an integral affine manifold
whereby the translational components of the transition maps are restricted to
also be integers. These structures are also interesting in their own right, for
example they have applications to physics \cite{HAMILTON2026105745}. One could
attempt to classify Lagrangian fibrations up to Lagrangian equivalence which
induce a given integral-integral affine structure on the base.

\subsection*{Investigate further the role of unipotence for integral affine manifolds}
A far reaching goal is to show the Markus conjecture holds for the case of
integral affine manifolds. As we have seen throughout this thesis, integral
affine manifolds have some special geometric features which may be possible to
exploit to achieve this. In Chapter \ref{ch:nilpotent-hol} we required a lot of
technical machinery in order to show that for affine manifolds with nilpotent
holonomy, unipotence of the linear holonomy representation is equivalent to
completeness of the affine structure.

However, with just the assumption of nilpotency, one can quickly see that at
least up to topological type, the Lagrangian fibrations over a fixed integral
affine structure depend only on the unipotent part of the representation. Let
$E_U$ denote the Fitting submodule, i.e. the unipotent part of the
representation. Fix an integral affine manifold $B$ and let $\Lambda$ denote
its corresponding period lattice. From the following short exact sequence of
$\pi_1(B)$-modules

$$
    0 \to E_U \to E \to E/E_U \to 0
$$

after taking the long exact sequence in cohomology and using that
$H^i(\pi_1(M), E/E_U) = 0$ by Lemma \ref{lemma:fitting-submodule} and Lemma
\ref{lemma:nilpotent-group-coho} as $\pi_1(M)$ is nilpotent, we get that
$H^2(\pi_1(M), E_U) \cong H^2(\pi_1(M), E)$. Hence if $M$ has nilpotent
holonomy the group corresponding to the Chern classes depends only on the
unipotent part of the representation, even without the assumptions of
compactness, completeness or parallel volume.

One might wonder if, for compact integral affine manifolds, this can be shown
to always hold. This is of course true if the Markus conjecture holds, but this
statement in itself is slightly weaker as we are not requiring the holonomy
representation itself to be unipotent which would imply completeness. Moreover,
inside the group $H^2(\pi_1(M), E)$ there is the subgroup of realisable Chern
classes. From the long exact sequence in cohomology of the above short exact
sequence, we see that $H^2(\pi_1(M), E_U) \subseteq H^2(\pi_1(M), E)$. It would
be interesting to determine whether the group of realisable Chern classes
always sits inside this group, which could hold even without the isomorphism
$H^2(\pi_1(M), E_U) \cong H^2(\pi_1(M), E)$.
